\subsection{dS Vacua and the Swampland --- arXiv:1901.02022}

\textbf{Authors:} Renata Kallosh, Andrei Linde, Evan McDonough, Marco Scalisi

\paragraph{主な主張}
uplift時のbackreactionによるポテンシャルの平坦化は、四次元KKLT有効理論の妥当な領域でde Sitter真空の形成を一般には妨げないと論じる\cite{Kallosh:2019axr}。

\paragraph{新規性}
10次元の反論の前提を再検討し、racetrack型のKL構成では小さなupliftが可能で、弱い重力予想と両立するパラメータ領域もあることを示す。

\paragraph{位置付け}
de Sitter構成への批判に対し、超重力有効理論の制御とbackreactionの大きさを区別して応答する研究。

\paragraph{コメント（読書メモ）}
de Sitter構成を支持する立場の論文として読んだ。ランドスケープの課題は、de Sitter真空の実現機構と宇宙項問題の解決の二つである。著者らは、将来$w\neq-1$が観測されてもno-de-Sitter予想の証明にはならないと論じる。その理由として、当時対象とされた暗黒エネルギー模型が観測と$4.5\sigma$以上の緊張にあるという評価と、de Sitter真空を含むランドスケープでも$w\neq-1$の模型を構成できる点を挙げる。さらに、漸近的なde Sitter状態へ転がり落ちる可能性がquintessenceの構成を簡単にし得る、と述べている。
